\scanpage{40}

On en conclut un isom canonique
\diagram{\boxed{\begin{tikzcd}[ampersand replacement=\&, column sep=large]
H^1(X,A) \& \operatorname{Ext}^1(A',\underline{\operatorname{Pic}}_X) \arrow[l,"\sim"']
\end{tikzcd}}}

\margin{Cet isomorphisme reste-t-il valable pour $X$ singulier ??? Il est possible qu'on ait besoin de $H^i(X,G_m)=0$ pour $i\unsure{\geq|>}2$, \unsure{vrai seulement si} $X$ non sing. ??}

($A$, variété abélienne), qui se compare à la formule de dualité
\[
\boxed{H^1(X,G)\simeq\operatorname{Hom}(D(G),\underline{\operatorname{Pic}}_X)}
\]
pour $G$ commutatif affine.

Remarque, On a des hom canoniques
\diagram{\begin{tikzcd}[ampersand replacement=\&, column sep=large,row sep=large]
\operatorname{Ext}^1(\underline{\operatorname{Alb}}_X,A) \arrow[r] \arrow[d,"\unsure{\wr}"] \& H^1(X,A)\\
\operatorname{Ext}^1(A',\underline{\operatorname{Alb}}_X') \arrow[d,"\unsure{\wr}"] \&\\
\operatorname{Ext}^1(A',\underline{\operatorname{Pic}}^\tau) \&
\end{tikzcd}}

dont on constate que c'est des \uline{isomorphismes} \uline{\unsure{donc|tous}}, qui définissent sur les \unsure{deux|trois} groupes en question l'isomorphisme envisagé plus haut.
\margin{39}
